%% =====================================================================
%%  1 / 7   Introduction  引言                                   成员 1
%% =====================================================================
\section{Introduction}
\label{sec:intro}

\paragraph{Background.}
Clinical artificial intelligence (AI) has spent a decade being useful in the passive voice. Systems read images, flag physiological deterioration and rank differential diagnoses, but they return a judgement and stop. That division of labour was stable enough to be codified in regulation and in training: the clinician decided, the software advised, and the interface between them was a screen. Two developments have disturbed it. The first is capability. Large language models have been wrapped into agents that plan, call external tools, write to records and execute multi-step procedures, so that the output of a clinical system is increasingly not a recommendation but a completed action \citep{wang2025baymax}. The second is institutional. Such systems are being deployed into workflows whose volume makes case-by-case review impractical, and they are marketed on the promise that review is no longer needed.

\paragraph{Problem.}
Evaluation has not changed to match. Benchmarks measure task success, and success is reported as if it were the only property that matters. On \textsc{MedAgentBench}, the most widely used virtual electronic health record environment, the strongest model completes 69.67\% of tasks \citep{jiang2025medagentbench}. A system that fails roughly one task in three is a legitimate object of study; it is not an obvious object of delegation. Yet the literature reports the number and leaves the delegation unexamined.

Two research communities describe the same phenomenon without meeting. Engineering work asks how to build agents that act correctly, and improves architectures, memory and tool use. Human-factors work asks what happens to people who work alongside automation, and has documented complacency, skill decay and misplaced mistrust for four decades \citep{parasuraman2010}. Neither asks the question that now matters: under what conditions may an agent be authorized to act at all?

A terminological looseness compounds the problem. ``Medical agent'' is used for systems ranging from a single-turn question answerer to an autonomous ward-round simulator. Without a stable referent, claims about autonomy cannot be compared, and a gain measured at one level is read as evidence about another.

\paragraph{Contributions.}
This paper makes four contributions.
\begin{itemize}
  \item We give a scoping definition of a medical agent --- planning, tool use and multi-step execution, as distinct from single-turn question answering --- and use it to bound the corpus.
  \item We conduct a purposive review of seventeen studies across six databases, placing engineering and human-factors sources in one corpus and coding each for its locus of action and level of autonomy.
  \item We identify four themes --- deskilling, the supervision illusion, accountability diffusion, and trust under uncertainty --- and show that all four are symptoms of one missing object: an explicit authorization boundary.
  \item We propose \emph{governed delegation}, which specifies four requirements for any agent permitted to act, and set out a $2\times2$ vignette protocol that tests the advisor-to-agent shift directly.
\end{itemize}

\paragraph{Roadmap.}
Section~\ref{sec:litreview} reviews the three literatures and states the gap. Section~\ref{sec:method} describes the review protocol and the experimental design. Section~\ref{sec:results} reports the coding and the four themes. Section~\ref{sec:discussion} develops governed delegation and answers two objections, and Section~\ref{sec:conclusion} concludes.
